\chapter{Dérivés de change}\label{ch:dv:fx-derivatives}

Le 11 août 2015, la banque centrale de Chine a modifié la façon dont elle fixait la parité centrale quotidienne du renminbi face au dollar.
La parité suivrait désormais le cours de clôture de la veille, et non plus un objectif fixé d'avance. Dans les deux jours qui
suivirent, le renminbi a perdu 2.8~\% face au dollar, et les devises asiatiques ont suivi le mouvement. Pour un client détenant un
\omterm{def:dv:fx-derivatives:tarf}{target redemption forward} (TARF), une structure qui rapporte peu tant que le taux reste d'un côté d'un strike et perd deux fois
plus vite de l'autre, un mouvement de cette ampleur fait la différence entre un produit qui se rembourse tôt avec un petit gain et un
produit qui arrive à échéance avec une lourde perte. Les options de change se cotent de manière très standardisée, en volatilité par delta, et leurs exotiques sont
évaluées à partir de ces cotations. Ce chapitre revient sur les conventions de valorisation, construit la
règle d'ajustement du smile propre au marché, vanna--volga, valorise les barrières et les \omterm{def:dv:fx-derivatives:tarf}{target redemption forwards} sous des modèles qui calent
le même smile, et présente le modèle qui combine \omterm{def:dv:local-volatility:model}{volatilité locale} et volatilité stochastique.

\section{Conventions de valorisation revisitées}

Le livre~2, chapitre~19, exposait comment les options de change sont cotées~: en volatilité, par delta, avec le
straddle à la monnaie, le risk reversal et le butterfly comme les trois instruments cotés pour chaque échéance. Le modèle sous-jacent aux
cotations est le suivant.

\begin{definition}[Modèle de Garman--Kohlhagen]\label{def:dv:fx-derivatives:gk}
Le \emph{modèle de Garman--Kohlhagen}\index{modèle de Garman--Kohlhagen} est le \omterm{def:dv:black-scholes-three-ways:model}{modèle de Black--Scholes} appliqué à un taux de change $S$ (unités de devise de cotation
par unité de devise de base)~: la devise de base est un actif versant un rendement continu égal à son taux d'intérêt $r_f$,
le taux de la devise de cotation $r_d$ sert à l'actualisation, et un call sur la devise de base vaut
$e^{-r_fT}S\,\Phi(d_1)-e^{-r_dT}K\,\Phi(d_2)$ avec $d_{1,2}=\bigl(\ln(F/K)\pm\frac12\sigma^2T\bigr)/(\sigma\sqrt T)$ et
$F=Se^{(r_d-r_f)T}$.
\end{definition}

Trois conventions que le modèle laisse ouvertes déterminent chaque cotation~: quel delta (spot ou forward, ajusté ou non de la prime),
quel point à la monnaie retenir (le straddle delta neutre pour la plupart des paires), et comment le butterfly est défini. Le composant
\texttt{firm.fxsmile} du livre~2 les implémente. Ce chapitre retient une paire dont le marché est connu exactement, afin que chaque méthode puisse être
vérifiée par rapport à la vérité. La paire oppose le dollar à une devise émergente, avec un spot de 7.00, un taux de la devise de cotation de 3~\% et un taux du dollar
de 5~\%, et son marché est un \omterm{def:dv:stochastic-volatility:heston}{modèle de Heston} (chapitre~10) avec $v_0=0.0036$, $\kappa=1.5$, $\bar v=0.0049$, $\eta=0.30$ et $\rho=+0.5$. La
corrélation est positive car, dans ces paires, les calls sur le dollar, c'est-à-dire la protection contre la baisse de la devise émergente, sont le
côté le plus cher.

\begin{example}[L'écran du courtier]\label{ex:dv:fx-derivatives:quotes}
À un an, le marché implique une volatilité à la monnaie de 5.32~\% (strike du straddle delta neutre 6.871, contre un forward de 6.861). Le
risk reversal à 25 delta vaut $+1.89$ points de volatilité et le butterfly $+0.46$, avec des strikes à 25 delta en 6.662 et 7.177. À 10 delta
le risk reversal vaut $+3.77$ et le butterfly $+1.71$ (strikes 6.442 et 7.702).
\end{example}

\section{La méthode vanna--volga}

Les dealers d'options à barrière et d'options touch avaient besoin d'une règle qui les valorise à partir des trois instruments cotés, sans
modèle complet de smile. La règle qu'ils ont adoptée couvre le risque de volatilité de l'exotique avec les
trois instruments et en facture le coût.

\begin{definition}[Méthode vanna--volga]\label{def:dv:fx-derivatives:vv}
La \emph{méthode vanna--volga}\index{méthode vanna--volga} valorise une option comme son prix de Black à la volatilité à la monnaie, plus
le coût de smile d'un portefeuille des trois instruments cotés (put à 25 delta, à la monnaie, call à 25 delta) qui reproduit
les sensibilités \omterm{def:dv:greeks-and-the-hedging-pnl:greeks}{vega}, \omterm{def:dv:greeks-and-the-hedging-pnl:greeks}{vanna} et \omterm{def:dv:greeks-and-the-hedging-pnl:greeks}{volga} de l'option (chapitre~4), toutes calculées à la volatilité à la monnaie.
\end{definition}

Pour une vanille de strike $K$, les poids $w_i$ résolvent un système linéaire trois par trois en \omterm{def:dv:greeks-and-the-hedging-pnl:greeks}{vega}, \omterm{def:dv:greeks-and-the-hedging-pnl:greeks}{vanna} et \omterm{def:dv:greeks-and-the-hedging-pnl:greeks}{volga} des piliers, et le prix
vaut $C_{\mathrm{BS}}(K,\sigma_{\mathrm{ATM}})+\sum_iw_i\bigl(C^{\mathrm{mkt}}_i-C^{\mathrm{BS}}_i\bigr)$. Elle reproduit les trois cotations
exactement, et entre elles elle donne un smile lisse, ce qui explique que la méthode soit très employée sur les marchés de change pour construire un smile complet
à partir de trois cotations. Son hypothèse sous-jacente a été décrite comme une volatilité implicite «~plate mais stochastique~». En dehors des
piliers, elle extrapole, et c'est là qu'elle peut se tromper.

\begin{example}[Vanna--volga face au vrai smile]\label{ex:dv:fx-derivatives:vvsmile}
Construit à partir des trois piliers à 25 delta de l'\cref{ex:dv:fx-derivatives:quotes}, le smile vanna--volga reproduit le marché à
quelques centièmes de point près entre les piliers. Au strike du call à 10 delta, il donne 7.94~\% contre 8.91~\% pour le marché, soit un point de moins
(\cref{fig:dv:fx-derivatives:smile}). Les cotations à 10 delta portent une information que les trois piliers ne portent pas.
\end{example}

\begin{omfigure}
\begin{tikzpicture}
\begin{axis}[width=10.5cm,height=5cm,
  xlabel={strike},ylabel={volatilité implicite à un an (\%)},
  tick label style={font=\small},label style={font=\small},
  xmin=6.35,xmax=7.8,ymin=4,ymax=9.8,grid=major,grid style={gray!25},
  legend columns=3,legend cell align=left,legend style={font=\footnotesize,at={(0.5,-0.3)},anchor=north,draw=none,fill=none,/tikz/every even column/.append style={column sep=8pt}},
  table/col sep=comma]
\addplot[gray,very thick] table[x=k,y=market]{figdata/derivatives/20-fx-derivatives/smile.csv};
\addplot[omThm,very thick,dashed] table[x=k,y=vv]{figdata/derivatives/20-fx-derivatives/smile.csv};
\addplot[black,only marks,mark=*,mark size=2pt] table[x=k,y=vol]{figdata/derivatives/20-fx-derivatives/pillars.csv};
\legend{marché (Heston),vanna--volga,les trois piliers}
\end{axis}
\end{tikzpicture}

\omcaption{Smile à un an de la paire du chapitre, du put à 10 delta au call à 10 delta~: le marché, et le smile vanna--volga construit à partir
du put à 25 delta, du straddle à la monnaie et du call à 25 delta. Exact aux piliers, proche entre eux, bas dans l'aile des calls. Données~: le
tutoriel.}\label{fig:dv:fx-derivatives:smile}
\end{omfigure}

Pour un touch ou une barrière, la même idée s'applique au \omterm{def:dv:greeks-and-the-hedging-pnl:greeks}{vanna} et au \omterm{def:dv:greeks-and-the-hedging-pnl:greeks}{volga} propres de l'exotique. Les coûts sont ceux du risk reversal (le
prix de marché du \omterm{def:dv:greeks-and-the-hedging-pnl:greeks}{vanna}) et du butterfly (le prix de marché du \omterm{def:dv:greeks-and-the-hedging-pnl:greeks}{volga}). La pratique de marché pondère ensuite la correction par une probabilité,
généralement la probabilité de ne pas toucher la barrière, car une barrière touchée tôt cesse d'exiger une couverture de smile. Cette pondération est
une convention et non un résultat démontré, et les desks en retiennent des variantes différentes. La section suivante mesure ce qu'elle produit.

\section{Barrières et target redemption forwards}

\begin{example}[One-touches évalués de cinq façons]\label{ex:dv:fx-derivatives:onetouch}
Une one-touch d'un an versant 1 à l'échéance si le taux touche une barrière haute, avec suivi continu (\cref{fig:dv:fx-derivatives:onetouch}).
Pour une barrière à 7.20, Black à la volatilité à la monnaie donne 45.5~\% du paiement, vanna--volga 34.5~\%, \omterm{def:dv:local-volatility:model}{volatilité locale} 43.3~\%,
\omterm{def:dv:local-volatility:model}{volatilité locale} stochastique 42.2~\% et Heston 38.9~\%. À 7.60, ils donnent 6.0~\%, 21.8~\%, 14.3~\%, 14.0~\% et 13.3~\%. À 8.00, ils donnent 0.4~\%,
6.7~\%, 6.5~\%, 6.4~\% et 6.2~\%. Les trois modèles calent le même smile et diffèrent jusqu'à 4.4 points. Le prix à la monnaie est très éloigné,
et la convention vanna--volga surestime les modèles pour les barrières intermédiaires de plus de moitié.
\end{example}

\begin{omfigure}
\begin{tikzpicture}
\begin{axis}[width=10.5cm,height=5.4cm,
  xlabel={barrière haute},ylabel={prix de la one-touch (\% du paiement)},
  tick label style={font=\small},label style={font=\small},
  xmin=7.15,xmax=8.05,ymin=0,ymax=48,grid=major,grid style={gray!25},
  legend columns=3,legend cell align=left,legend style={font=\footnotesize,at={(0.5,-0.28)},anchor=north,draw=none,fill=none,/tikz/every even column/.append style={column sep=8pt}},
  table/col sep=comma]
\addplot[gray,thick,dashed,mark=o,mark size=1.6pt] table[x=h,y=bs]{figdata/derivatives/20-fx-derivatives/one_touch.csv};
\addplot[omProp,thick,mark=triangle*,mark size=1.8pt] table[x=h,y=vv]{figdata/derivatives/20-fx-derivatives/one_touch.csv};
\addplot[omDef,very thick,mark=square*,mark size=1.6pt] table[x=h,y=lv]{figdata/derivatives/20-fx-derivatives/one_touch.csv};
\addplot[black,thick,mark=diamond*,mark size=1.8pt] table[x=h,y=slv]{figdata/derivatives/20-fx-derivatives/one_touch.csv};
\addplot[omThm,very thick,mark=*,mark size=1.6pt] table[x=h,y=sv]{figdata/derivatives/20-fx-derivatives/one_touch.csv};
\legend{Black à la monnaie,vanna--volga,volatilité locale,{volatilité locale stochastique, mélange 0.5},Heston}
\end{axis}
\end{tikzpicture}

\omcaption{One-touches supérieurs à un an sur la paire du chapitre (spot 7.00), par barrière. Les trois modèles ajustent le même smile~; la volatilité de Black
à la monnaie l'ignore~; vanna--volga applique la règle empirique du marché. Données~: le tutoriel.}\label{fig:dv:fx-derivatives:onetouch}
\end{omfigure}

Les désaccords entre modèles relèvent à nouveau de la dynamique des chapitres 9 à 11. La \omterm{def:dv:local-volatility:model}{volatilité locale} fait du smile une fonction du spot, si bien que lorsque le taux
monte vers la barrière, il entre dans l'aile des calls à forte volatilité et touche plus souvent. Heston conserve une part de la volatilité dans un
processus de variance, corrélé au taux mais non fonction de celui-ci. Une touche dépend de toute la trajectoire, donc de ces dynamiques, et
non seulement du smile terminal.

Les \omterm{def:dv:the-structured-products-business:product}{produits structurés} les plus notoires du marché des changes reposent sur des forwards plutôt que sur des options.

\begin{definition}[Target redemption forward]\label{def:dv:fx-derivatives:tarf}
Un \emph{target redemption forward}\index{target redemption forward} (TARF) est une suite de contrats forward à un strike amélioré, un
par date de fixing, qui prend fin dès que les gains cumulés du client atteignent un montant cible. Les fixings situés du côté défavorable du strike sont
réglés sur un notionnel à effet de levier (généralement le double), sans limite sur les pertes.
\end{definition}

\begin{definition}[Accumulator]\label{def:dv:fx-derivatives:accumulator}
Un \emph{accumulator}\index{accumulator} est l'équivalent sur actions ou sur devises dans lequel le client achète (ou vend) une quantité fixe à un
strike décoté à chaque fixing tant qu'une barrière à extinction n'a pas été touchée, et le double de la quantité lorsque le cours est du mauvais
côté du strike.
\end{definition}

Ces deux produits offrent au client un meilleur taux que le forward, en contrepartie d'une position vendeuse d'options avec effet de levier. La cible ou la barrière à extinction
plafonne rapidement le gain potentiel, et l'effet de levier laisse le risque de baisse ouvert. Le client a vendu de la volatilité, et souvent davantage qu'il ne l'imagine.
Le produit est valorisé à coût nul en choisissant le strike.

\begin{example}[Un TARF sur la paire]\label{ex:dv:fx-derivatives:tarf}
Le client vend un million de dollars à chacune des douze dates de fixing mensuelles, au strike $K$~: un gain de $K-S$ par dollar lorsque le fixing est inférieur à
$K$, une perte de $2(S-K)$ lorsqu'il est au-dessus, et extinction dès que les gains atteignent 0.30 par dollar. Dans le \omterm{def:dv:fx-derivatives:slv}{modèle à volatilité locale stochastique} (poids de mélange 0.5), le
strike à coût nul est de 7.312, contre un spot de 7.00 et un forward à un an de 6.861. Le client se voit offrir un taux supérieur de 6.6~\% à celui du
forward. Le TARF s'éteint dans 98.1~\% des trajectoires, après 1.73 fixings en moyenne, et 95~\% des trajectoires se terminent avec un gain d'au moins 126\,000
(la \cref{fig:dv:fx-derivatives:tarf}, à gauche). La valeur nulle provient des 5~\% restants. Si le taux augmente de 10~\% au cours du premier trimestre, la
perte espérée du client est de 7.12 millions par million de notionnel mensuel~: 1.01 million au cours du premier trimestre et le reste sur neuf
fixings supplémentaires au double du notionnel. Une hausse de 5~\% coûte 0.99 million, tandis qu'une hausse allant jusqu'à 2.5~\% permet encore au TARF de s'éteindre avec le
plein gain de 0.30 million (la \cref{fig:dv:fx-derivatives:tarf}, à droite).
\end{example}

\begin{omfigure}
\begin{tikzpicture}
\begin{groupplot}[group style={group size=2 by 1,horizontal sep=1.6cm},
  width=6.6cm,height=5cm,
  tick label style={font=\small},label style={font=\small},grid=major,grid style={gray!25},table/col sep=comma]
\nextgroupplot[ybar,bar width=6pt,xlabel={fixings jusqu'à l'extinction},ylabel={part des trajectoires (\%)},xmin=0.3,xmax=12.7,ymin=0,ymax=72,
  xtick={1,2,4,6,8,10,12},title={\small durée de vie}]
\addplot[omThm,fill=omThm!40] table[x=n,y=share]{figdata/derivatives/20-fx-derivatives/tarf_fixings.csv};
\nextgroupplot[xlabel={variation du taux au cours du premier trimestre (\%)},ylabel={P\&L espéré du client (million)},xmin=-6,xmax=11,ymin=-8,ymax=1,
  title={\small une variation défavorable pour le client}]
\addplot[omDef,very thick,mark=*,mark size=1.6pt] table[x=move,y=pnl]{figdata/derivatives/20-fx-derivatives/tarf_scenarios.csv};
\draw[gray] (axis cs:-6,0) -- (axis cs:11,0);
\end{groupplot}
\end{tikzpicture}

\omcaption{Un TARF à coût nul sur la paire du chapitre (le client vend le dollar à 7.312, avec le double du notionnel au-delà de ce niveau, cible 0.30). À gauche~: le
nombre de fixings avant extinction~: la plupart s'éteignent au premier ou au deuxième. À droite~: le P\&L espéré du client, par million de notionnel mensuel,
si le taux varie d'une amplitude donnée au cours du premier trimestre~; les fixings d'un TARF qui ne parvient pas à s'éteindre sont réglés au double du
notionnel. Données~: le tutoriel.}\label{fig:dv:fx-derivatives:tarf}
\end{omfigure}

Le mouvement d'août 2015 montre à quel point le produit est dangereux. Pour un client détenant un TARF de ce type, une dépréciation durable de la
devise émergente est exactement le scénario dont le strike à coût nul a été le prix. La cible n'est jamais atteinte, et les fixings à effet de levier s'accumulent. Du
côté du dealer, la situation n'est pas plus simple. Le delta d'un TARF s'inverse à mesure que le gain cumulé s'approche de la cible, et sa valeur dépend de la
dynamique du smile au fil de la trajectoire, comme pour les touches ci-dessus.

\section{Volatilité locale stochastique}

La \omterm{def:dv:local-volatility:model}{volatilité locale} ajuste le smile mais se trompe sur la dynamique dans un sens~; la volatilité stochastique la capture en partie et ne peut
pas ajuster exactement le smile. Le compromis naturel consiste à les combiner.

\begin{definition}[Modèle à volatilité locale stochastique]\label{def:dv:fx-derivatives:slv}
Un \emph{modèle à volatilité locale stochastique}\index{modèle à volatilité locale stochastique} (stochastic-local volatility model, SLV) définit la volatilité du sous-jacent comme le
produit d'une fonction du temps et du spot et d'un facteur stochastique~:
\[
\frac{dS_t}{S_t}=(r_d-r_f)\,dt+L(t,S_t)\sqrt{v_t}\,dW^1_t,\qquad dv_t=\kappa(\bar v-v_t)\,dt+\lambda\eta\sqrt{v_t}\,dW^2_t,
\]
avec un poids de mélange $\lambda\in[0,1]$ qui gouverne la volatilité de la variance~: $\lambda=0$ réduit le modèle à la \omterm{def:dv:local-volatility:model}{volatilité locale} et $\lambda=1$
avec $L(t,S)=1$, il se ramène au \omterm{def:dv:stochastic-volatility:heston}{modèle de Heston}.
\end{definition}

\begin{definition}[Fonction de levier]\label{def:dv:fx-derivatives:leverage}
La \emph{fonction de levier}\index{fonction de levier} $L(t,S)$ d'un \omterm{def:dv:fx-derivatives:slv}{modèle à volatilité locale stochastique} est la fonction qui fait en sorte que le modèle
reproduise le prix de chaque option vanille. Par la \omterm{def:dv:local-volatility:projection}{projection markovienne} du chapitre 9, elle vérifie
\[
L(t,S)^2\,\E\bigl[v_t\mid S_t=S\bigr]=\sigma_{\mathrm{loc}}(t,S)^2,
\]
où $\sigma_{\mathrm{loc}}(t,S)$ est la \omterm{def:dv:local-volatility:model}{volatilité locale} de marché.
\end{definition}

L'équation est circulaire~: $L(t,S)$ dépend de l'espérance conditionnelle, qui dépend de la loi des trajectoires, qui dépend de $L(t,S)$.
Guyon et Henry-Labordère ont montré que la méthode particulaire de McKean la résout exactement. Simuler un vaste ensemble de particules en avançant dans le temps.
À chaque pas, estimer $\E[v_t\mid S_t=S]$ à partir des particules elles-mêmes, par exemple en les répartissant en classes de $S$, puis fixer
$L(t,S)$ à partir de l'équation avant de passer au pas suivant. La calibration coûte une seule simulation. La construction procède avec 50\,000 particules
et 40 classes.

\begin{example}[La fonction de levier de la paire]\label{ex:dv:fx-derivatives:leverage}
Avec un mélange $\lambda=0.5$, la $L(t,S)$ calibrée au voisinage du forward vaut 0.76 à trois mois, 0.73 à six mois et 0.79 à un an. Elle
monte à environ 1.4 ou 1.5 sur les deux ailes (la \cref{fig:dv:fx-derivatives:leverage}). Les trois modèles reproduisent le smile à un an~: le put 10-delta,
le put 25-delta, le straddle à la monnaie, le call 25-delta et le call 10-delta, la \omterm{def:dv:local-volatility:model}{volatilité locale} stochastique donne 5.22~\%, 4.90~\%,
5.38~\%, 6.77~\% et 8.96~\%, contre 5.14~\%, 4.84~\%, 5.32~\%, 6.72~\% et 8.91~\% sur le marché. Les écarts, inférieurs à un dixième de point, proviennent
des pas d'Euler quotidiens de la simulation.
\end{example}

\begin{omfigure}
\begin{tikzpicture}
\begin{axis}[width=10.5cm,height=5cm,
  xlabel={spot},ylabel={fonction de levier $L(t,S)$},
  tick label style={font=\small},label style={font=\small},
  xmin=6.0,xmax=8.6,ymin=0.6,ymax=1.6,grid=major,grid style={gray!25},
  legend columns=3,legend cell align=left,legend style={font=\footnotesize,at={(0.5,-0.3)},anchor=north,draw=none,fill=none,/tikz/every even column/.append style={column sep=8pt}},
  table/col sep=comma]
\addplot[omProp,thick] table[x=s3m,y=l3m]{figdata/derivatives/20-fx-derivatives/leverage_axes.csv};
\addplot[omDef,very thick,dashed] table[x=s6m,y=l6m]{figdata/derivatives/20-fx-derivatives/leverage_axes.csv};
\addplot[omThm,very thick] table[x=s12m,y=l12m]{figdata/derivatives/20-fx-derivatives/leverage_axes.csv};
\legend{3 mois,6 mois,1 an}
\end{axis}
\end{tikzpicture}

\omcaption{La \omterm{def:dv:fx-derivatives:leverage}{fonction de levier} du \omterm{def:dv:fx-derivatives:slv}{modèle à volatilité locale stochastique} (mélange 0.5) calibré par la méthode particulaire sur la paire du chapitre,
à trois dates, sur l'intervalle couvert par les particules. Inférieure à un près du forward, supérieure à un dans les ailes~: la composante locale fournit le smile que
la part en volatilité de variance, de poids 0.5, laisse de côté. Données~: le tutoriel.}\label{fig:dv:fx-derivatives:leverage}
\end{omfigure}

Le poids de mélange est le paramètre libre du modèle. Il n'est pas ajusté aux vanilles, que chaque $\lambda$ reproduit à l'identique. Il se choisit à partir des
produits qui dépendent de la dynamique, des touches et des barrières avant tout, en le calant sur leurs prix lorsque le marché les cote, ou sur
le \omterm{def:dv:local-volatility:forward}{smile forward} lorsqu'il ne les cote pas. Dans l'exemple, passer de la \omterm{def:dv:local-volatility:model}{volatilité locale} au \omterm{def:dv:stochastic-volatility:heston}{modèle de Heston} fait baisser la one-touch à 7.20 de 43.3~\% à 38.9~\%.
La \omterm{def:dv:local-volatility:model}{volatilité locale} stochastique avec $\lambda=0.5$ se situe entre les deux, à 42.2~\%.

\section{Un portefeuille d'exotiques de change}

Sur des dizaines de paires, un desk d'exotiques de change gère des milliers de touches, de barrières, de TARF et d'\omterm{def:dv:fx-derivatives:accumulator}{accumulators}. Ses questions quotidiennes découlent
de ce chapitre.
\begin{itemize}
\item \textbf{Quel modèle, quel mélange.} Un modèle stochastique-local par paire, son poids de mélange calibré sur le marché liquide des touches, et
le prix \omterm{def:dv:greeks-and-the-hedging-pnl:greeks}{vanna-volga} conservé en parallèle comme référence du marché pour la cotation.
\item \textbf{Où se situe le risque.} Le \omterm{def:dv:greeks-and-the-hedging-pnl:greeks}{vega} par échéance, ainsi que le \omterm{def:dv:greeks-and-the-hedging-pnl:greeks}{vanna} et le \omterm{def:dv:greeks-and-the-hedging-pnl:greeks}{volga}, que concentre un portefeuille de barrières près de ses barrières. Les portefeuilles de TARF ajoutent
une exposition importante à la dérive du taux et aux sauts, puisqu'un mouvement soutenu est précisément ce contre quoi ils sont vendeurs.
\item \textbf{Ce que le modèle ne voit pas.} Les devises gérées évoluent par paliers lorsque la politique change, comme l'a montré le changement de fixing d'août 2015.
Les 2.8~\% en deux jours constituaient un saut pour un modèle calibré sur une paire qui avait évolué près du côté faible d'une bande étroite. Le risque de saut et le
régime de politique relèvent des réserves (chapitre 27), et non de la volatilité.
\end{itemize}

\section{Tutoriel~: de trois cotations à un TARF}\label{tut:dv:fx-derivatives:tut}

\begin{tutorial}
\textbf{Objectif.} Lire les cotations des courtiers d'un marché connu, construire le smile \omterm{def:dv:greeks-and-the-hedging-pnl:greeks}{vanna-volga} et la one-touch, calibrer un \omterm{def:dv:fx-derivatives:slv}{modèle à volatilité locale stochastique}
par la méthode particulaire, et pricer un TARF. \textbf{État final~:} les quatre figures et les nombres du problème du week-end.
\begin{tutsteps}
\item \textbf{Poids \omterm{def:dv:greeks-and-the-hedging-pnl:greeks}{vanna-volga} et prix}~:
\omcode{code/firm/fxvol/firm_fxvol.py}{30}{44}{Poids vanna-volga et prix de la vanille.}\label{lst:dv:fx-derivatives:vv}
\item \textbf{La méthode particulaire}~: regrouper les particules par classes de spot, estimer la variance conditionnelle, fixer le levier et le pas~:
\omcode{code/firm/fxvol/firm_fxvol.py}{108}{139}{Fonction de levier par la méthode particulaire.}\label{lst:dv:fx-derivatives:particle}
\item \textbf{Le TARF sur fixings simulés}~:
\omcode{code/firm/fxvol/firm_fxvol.py}{171}{192}{Flux de trésorerie du target redemption forward.}\label{lst:dv:fx-derivatives:tarf}
\item \textbf{Exécuter} \texttt{dv\_fx.quotes()}, \texttt{vv\_smile()}, \texttt{one\_touches()}, \texttt{tarf\_summary()}, \texttt{tarf\_scenario(0.10)} et
\texttt{fig\_fx.py}. La \omterm{def:dv:local-volatility:model}{volatilité locale} de marché est celle de Dupire, sur une grille de prix de Heston (une évaluation de Fourier par échéance).
\end{tutsteps}
\textbf{Pour aller plus loin.} Calibrer le poids de mélange pour que la one-touch du modèle à 7.60 corresponde au prix \omterm{def:dv:greeks-and-the-hedging-pnl:greeks}{vanna-volga}~; pricer le TARF en \omterm{def:dv:local-volatility:model}{volatilité locale}
puis en Heston et comparer les strikes à coût nul~; ajouter un knock-out au TARF.
\end{tutorial}

\section{Construction~: volatilité de change}\label{bld:dv:fx-derivatives:fxvol}

\begin{build}
\bfield{Rôle} Le moteur d'exotiques de change de la firme miniature~: la convention de marché pour les touches et les barrières, le modèle stochastique-local avec lequel le
desk valorise, et les flux de trésorerie du TARF et de l'\omterm{def:dv:fx-derivatives:accumulator}{accumulator} que son desk de \omterm{def:dv:the-structured-products-business:product}{produits structurés} vend.
\bfield{Interface} \texttt{vv\_weights}, \texttt{vv\_price}, \texttt{vv\_implied}~; \texttt{up\_hit\_probability}, \texttt{one\_touch\_bs},
\texttt{one\_touch\_vv}~; \texttt{SLV(v0, kappa, vbar, eta, rho, mix)}~; \texttt{calibrate\_leverage(model, local\_vol, s0, rd, rf, t\_end,
\ldots, n, bins)}~; \texttt{simulate\_slv(model, leverage, s0, rd, rf, t\_end, n, \ldots)}~; \texttt{tarf\_client\_pnl(fixings, strike, target,
leverage, notional)}.
\bfield{Règles} Les cotations sont converties selon les conventions du livre~2 avant que tout modèle ne les voie~; la \omterm{def:dv:fx-derivatives:leverage}{fonction de levier} est recalibrée à chaque
changement du smile~; le poids de mélange est un paramètre documenté et relu~; les touches issues de trajectoires quotidiennes sont décalées vers un suivi continu.
\bfield{Tests d'acceptation} \texttt{code/firm/fxvol/tests/}~: le \omterm{def:dv:greeks-and-the-hedging-pnl:greeks}{vanna-volga} reproduit ses piliers et un smile plat~; la one-touch de Black coïncide avec une
simulation~; sans volatilité de variance, le modèle stochastique-local se réduit à la \omterm{def:dv:local-volatility:model}{volatilité locale} et reproduit un smile plat~; flux de trésorerie du TARF sur des
trajectoires écrites à la main.
\bfield{Pour aller plus loin} Double no-touches en SLV~; un poids de mélange calibré sur les cotations de touches~; \omterm{def:dv:fx-derivatives:accumulator}{accumulators} avec knock-outs quotidiens~; le delta et le
\omterm{def:dv:greeks-and-the-hedging-pnl:greeks}{vega} du TARF par fixing.
\end{build}

\begin{omsources}
\item M.~B.~Garman et S.~W.~Kohlhagen, «~Foreign currency option values~», \emph{Journal of International Money and Finance} 2(3) (1983)
231--237.
\item J.~Guyon et P.~Henry-Labordère, «~The smile calibration problem solved~», SSRN 1885032 (2011).
\item A.~Perederiy, «~Vanna--volga method for normal volatilities~», SSRN 3262146 (2018)~; F.~Rolloos, «~On the flat but stochastic implied
volatility assumption in the vanna--volga model~», SSRN 4287718 (2022).
\item Bank for International Settlements, \emph{Quarterly Review} (septembre 2015), «~EME vulnerabilities take centre stage~».
\end{omsources}

\section{Exercices}

\begin{exercise}[$\star$]\label{exo:dv:fx-derivatives:1}
Montrer que le call de Garman--Kohlhagen est la formule de Black sur la forward $F=Se^{(r_d-r_f)T}$, actualisée au taux $r_d$.
\end{exercise}

\begin{exercise}[$\star$]\label{exo:dv:fx-derivatives:2}
Expliquer pourquoi la corrélation entre le taux et sa variance est positive pour un dollar contre une devise émergente, et ce
qu'elle fait au risk reversal~?
\end{exercise}

\begin{exercise}[$\star$]\label{exo:dv:fx-derivatives:3}
Le strike du TARF du client est 7.312 alors que la forward vaut 6.861. Qu'a payé le client pour un taux meilleur de 6.6~\% que la forward~?
\end{exercise}

\begin{exercise}[$\star\star$]\label{exo:dv:fx-derivatives:4}
Expliquer pourquoi le smile \omterm{def:dv:greeks-and-the-hedging-pnl:greeks}{vanna-volga} est exact aux trois piliers et pourquoi il peut être faux à 10 delta.
\end{exercise}

\begin{exercise}[$\star\star$]\label{exo:dv:fx-derivatives:5}
La \omterm{def:dv:local-volatility:model}{volatilité locale} donne le prix de one-touch le plus élevé à 7.20 parmi les trois modèles. Expliquer.
\end{exercise}

\begin{exercise}[$\star\star$]\label{exo:dv:fx-derivatives:6}
Montrer qu'avec $\lambda=0$ la \omterm{def:dv:fx-derivatives:leverage}{fonction de levier} vaut $\sigma_{\mathrm{loc}}(t,S)/\sqrt{v(t)}$, où $v(t)$ est la trajectoire déterministe de la variance.
\end{exercise}

\begin{exercise}[$\star\star\star$]\label{exo:dv:fx-derivatives:7}
\emph{Codage.} Pricer le TARF du chapitre en \omterm{def:dv:local-volatility:model}{volatilité locale} puis en Heston, et comparer les strikes à coût nul et le nombre espéré de
fixings.
\end{exercise}

\begin{exercise}[$\star\star\star$]\label{exo:dv:fx-derivatives:8}
\emph{Repérer l'erreur.} «~Notre TARF se rembourse dans 98~\% des scénarios, donc pour un client c'est presque une amélioration gratuite du taux forward.~»
\end{exercise}

\section{Problème~: le TARF}

\begin{problem}[{Problème du week-end --- combien de temps il vit et ce qu'il coûte}]\label{pb:dv:fx-derivatives:1}
Un client vend un million de dollars par mois pendant un an contre la devise émergente du chapitre, par l'intermédiaire d'un TARF~: le double du notionnel lorsque le fixing
est au-dessus du strike, remboursement lorsque le gain accumulé atteint 0.30 par dollar. Le dealer fait son pricing avec une \omterm{def:dv:local-volatility:model}{volatilité locale} stochastique (mélange 0.5).

\medskip\noindent\textbf{Partie I --- Le marché.}
\begin{enumerate}
\item Donner la volatilité à un an à la monnaie, le risk reversal et le butterfly à 25 delta.
\item Que dit le risk reversal positif sur les craintes du marché~?
\item Donner le forward à un an, et expliquer pourquoi il est inférieur au spot.
\item Dans quelle mesure le \omterm{def:dv:fx-derivatives:slv}{modèle à volatilité locale stochastique} reproduit-il le smile à un an~?
\item Donner la \omterm{def:dv:fx-derivatives:leverage}{fonction de levier} au voisinage du forward à un an.
\end{enumerate}

\medskip\noindent\textbf{Partie II --- Le TARF.}
\begin{enumerate}[resume]
\item Donner le strike à coût nul et son amélioration par rapport au forward.
\item Donner la part des trajectoires qui se remboursent et le nombre espéré de fixings.
\item Quel est le gain du client sur la trajectoire typique, et d'où provient la valeur nulle~?
\item Donner les strikes à coût nul sous la \omterm{def:dv:local-volatility:model}{volatilité locale} et sous le \omterm{def:dv:stochastic-volatility:heston}{modèle de Heston}.
\item Pourquoi une cible plus élevée abaisse-t-elle le strike que le dealer peut proposer~?
\end{enumerate}

\medskip\noindent\textbf{Partie III --- Le mouvement.}
\begin{enumerate}[resume]
\item Le taux monte de 10~\% au cours du premier trimestre. Donner la perte espérée du client, et la part de celle-ci subie au premier trimestre.
\item Donner la même chose pour une hausse de 5~\%.
\item Combien de fixings le TARF vit-il, en moyenne, après une hausse de 10~\%~?
\item À quoi ressemble la couverture du dealer lorsque le gain accumulé approche de la cible~?
\item Mettre en relation ce scénario avec le mouvement du renminbi en août 2015.
\end{enumerate}

\medskip\noindent\textbf{Partie IV --- Jugement.}
\begin{enumerate}[resume]
\item Que le client aurait-il dû voir avant de négocier~?
\item Quel risque de modèle compte le plus pour le dealer~?
\item Vendriez-vous ce produit à une entreprise sans revenu naturel en dollars~?
\item Énoncer le \emph{résultat nommé}~: le nombre espéré de fixings avant remboursement, et la perte espérée du client par million lorsque le taux
évolue de 10~\% contre lui au premier trimestre.
\item En une phrase~: que vend un client de TARF~?
\end{enumerate}
\end{problem}

\section{Questions d'entretien}

\begin{interviewq}[$\star$ \iqroles{trader}]\label{iq:dv:fx-derivatives:1}
Expliquez les cotations à la monnaie, le risk reversal et le butterfly, et comment vous les transformeriez en smile.
\end{interviewq}

\begin{interviewq}[$\star\star$ \iqroles{chercheur}]\label{iq:dv:fx-derivatives:2}
Dérivez le prix \omterm{def:dv:greeks-and-the-hedging-pnl:greeks}{vanna-volga} d'une vanille. Que suppose la méthode~?
\end{interviewq}

\begin{interviewq}[$\star\star$ \iqroles{chercheur, développeur}]\label{iq:dv:fx-derivatives:3}
Qu'est-ce qu'un \omterm{def:dv:fx-derivatives:slv}{modèle à volatilité locale stochastique}, et comment calibreriez-vous sa \omterm{def:dv:fx-derivatives:leverage}{fonction de levier}~?
\end{interviewq}

\begin{interviewq}[$\star\star$ \iqroles{trader, risque}]\label{iq:dv:fx-derivatives:4}
Deux modèles qui ajustent le même smile diffèrent de quatre points sur une one-touch. Comment choisissez-vous le prix~?
\end{interviewq}

\begin{interviewq}[$\star\star$ \iqroles{risque, banque}]\label{iq:dv:fx-derivatives:5}
Décrivez un TARF du point de vue du client. Qu'est-ce qui peut mal tourner pour le client et pour la banque~?
\end{interviewq}

\begin{interviewq}[$\star\star\star$ \iqroles{trader, risque}]\label{iq:dv:fx-derivatives:6}
Une monnaie gérée peut être dévaluée par palier. Comment cela modifie-t-il le pricing et la couverture de votre portefeuille de barrières et de TARF~?
\end{interviewq}
